\chapter{Курсовая работа\+: Разработка серверной программы}
\hypertarget{md_README}{}\label{md_README}\index{Курсовая работа: Разработка серверной программы@{Курсовая работа: Разработка серверной программы}}
\label{md_README_autotoc_md0}%
\Hypertarget{md_README_autotoc_md0}%
Программа представляет собой TCP-\/сервер на языке C++, который принимает подключения от клиентов, проводит аутентификацию по логину и паролю (с использованием хэш-\/функции SHA-\/1 и соли), а затем принимает бинарные данные (векторы) и возвращает сумму их элементов.

Сборка и запуск

make

В первом окне\+:

chmod +x server\+\_\+static

./server\+\_\+static -\/T double -\/H SHA1 -\/S server

Во втором окне\+:

./server\+\_\+double -\/H SHA1 -\/S server

Вывод справки\+: ./vcalc\+\_\+server -\/h, ./vcalc\+\_\+server --help, ./vcalc\+\_\+server

Тестирование рабочих параметров\+: ./vcalc\+\_\+server --run -\/p 4444

Своя база данных\+: Создать файл базы\+: echo "{}user\+:123"{} \texorpdfstring{$>$}{>} my\+\_\+base.\+conf Запустить\+:./vcalc\+\_\+server --run -\/d my\+\_\+base.\+conf

Свой лог-\/файл\+: ./vcalc\+\_\+server --run -\/l my\+\_\+log.\+txt комбинация параметров\+: ./vcalc\+\_\+server --run -\/p 33333 -\/d clients.\+conf 